\section{Introduction}
\label{sec:introduction}

Fertility is falling almost everywhere. In the United States, the period
fertility rate fell from about 2.1 children per woman in 2007 to 1.62 in
2023, and the world as a whole is likely below replacement for the first
time \citep{worldBankWDI,fernandezVillaverdeNorrickTerraIncognita2026}. The consequences of this decline are vast, ranging from social security impacts to a slowdown in growth. As a result, this phenomenon has become a central concern for policymakers. The causes of the decline are elusive,%
\footnote{\citet{kearneyPuzzleFallingUS2022} find no economic or policy
factor that accounts for the decline. \citet{kearneyLevineWhyFertilityLow2025}
argue for ``shifting priorities'' across cohorts, a change in norms around
work, parenting, gender roles and leisure, shaped in part by new technologies
such as the smartphone and social media.
\citet{doepkeEconomicsFertilityNew2023} argue that in rich countries fertility now depends on how easily parents can combine careers and family.} but housing is often at the center of the
discussion on how to raise fertility in the future, both in public debate and
in policy.\footnote{In a fourteen-country
survey, housing is the third most cited reason for having fewer children than
desired \citep{unfpaSWP2025}. Germany, Greece, Hungary, Italy, Poland and
Spain subsidize or guarantee home loans for families, with terms that improve
with the number of children
\citep{kfw300,greeceSpitiMou2,hungaryCSOKPlus,italyFondoPrimaCasa,polandFamilyRepayment,spainICOJovenes};
South Korea lends at subsidized rates to households with a newborn
\citep{koreaNewbornLoan}, and Singapore gives families priority for public
flats \citep{singaporeFPPS}. In the 2024 US presidential campaign, Kamala
Harris proposed up to \$25,000 in down-payment assistance for first-time
buyers \citep{harrisCampaign2024}, and JD Vance argued that young people are
delaying families ``because they can't afford a place to raise a family''
\citep{vanceNewsNation2024}.} The logic is intuitive. Housing is the largest item in household budgets, a third of all spending
\citep{blsConsumerExpenditures2024}.  Children increase the pressure on this already high bill: they require physical space, so becoming a parent raises a household's demand for housing. On the other hand, prospective parents are young, with low earnings and savings, so that the extra room is particularly expensive exactly when it could shift the fertility decision.

Many policy proposals seek to help young families buy a home, intervening on the financial constraints by lowering the cost of borrowing or the down payment. If saving for a down payment is what keeps young couples out of a `family-sized' home, especially where larger homes have to be bought,\footnote{In the United States, 43\% of
owned homes have seven or more rooms, against 7.5\% of rented homes (2023
American Housing Survey).} this could raise births. But the extent to which this is quantitatively relevant is unclear. Easier credit lowers the down payment, but not the cost of the extra space, which a family pays every year its children live at home: whether it rents or owns, and in bad years as well as good ones. Couples who get access to eased borrowing, then, may decide to buy a house, and perhaps have children earlier, without necessarily having more of them.

This paper asks how the cost of family space shapes fertility, and what
housing policy can do about it. I answer it with a quantitative life-cycle model in which households face idiosyncratic earnings risk and choose whether to have children, whether to rent or own, how much housing to demand and how much to save. I find that housing policy raises births when it lowers the cost of having children, through cheaper space or credit tied to births, rather than just easing financial constraints and the desire to own. Removing the down
payment for all households and expanding the housing supply raise births by
about the same amount on impact, 1.3\% and 1.4\%. The gain from credit is gone
within four years, while the gain from supply keeps growing, to 3.0\% by
2063. Credit available only to parents raises births by 5.1\% on impact and
by 16.6\% by 2063.

The reason is what each policy changes for a couple deciding whether to have a child. Lowering the down payment makes a home easier to buy, but the flow cost of children is the same. As it does not lower the burden of parenthood, and it does not change the demand for space, most of the households who use it buy the small homes they would otherwise have rented. As a comparison, increasing housing supply lowers prices and rents, and with them the cost of family space in the present and future. Offering the credit only to parents is very effective, shifting the value of parenthood relative to childlessness.

These findings matter for a growing empirical literature that links both
housing costs and access to credit with births
\citep{dettlingHousePricesBirth2014,lovenheimFamilyWealthShocks2013,hacamoBabiesMortgageDeregulation2021,cummingMonetaryPolicyBirth2024,dettlingKearney2026,vandoornikFazioRamadoraiSkrastinsHousingFertility2025}.
My model is consistent with a positive response of births when credit becomes
easier, and it complements this evidence in two ways. First, a structural general-equilibrium model helps in disentangling the forces at play. By modeling the households' choice of children, tenure, housing and savings together, the model can shed light on the interaction between housing and children in equilibrium, and it can evaluate policies that have not been tried, such as credit available only to parents. Second, it takes a `long view'. Empirical designs typically operate at short horizons and given prices. Some studies look at longer timeframes
\citep{dettlingKearney2026,vandoornikFazioRamadoraiSkrastinsHousingFertility2025},
but they are necessarily retrospective: it takes more than two decades for
the effect of a policy on completed fertility to show. The model instead allows me to project forward. Starting from the economy of 2023, I follow each policy for several decades, as cohorts age and prices adjust.

The paper builds on evidence on how housing changes at parenthood. In the Panel Study of Income
Dynamics, three to four years after their first birth parents live in 1.44
more rooms than two to three years before it, a 31\% increase, and they are
27 percentage points more likely to own their home. Moves for more space
become a third more likely after the birth, while moves for a better
neighborhood do not. The demand for space jumps when a child arrives, and
owning is the main route to a large home.

In the model, households live through a life cycle and face uninsurable
earnings risk. They save, decide whether to have children, whether to rent or
own, and how much space to live in; a child requires a minimum amount of space
while it lives at home. The largest homes are only for sale, buying requires
a down payment, and owners can borrow against their home. Housing is supplied
along an upward-sloping supply curve, so house prices rise with the number of
households, and fertility choices feed back into the price of space a
generation later.

I approximate the United States economy of
2007 with a steady state in which births are at replacement. I assign some
parameters from external evidence and choose the others so that this steady
state matches fertility, housing and wealth moments of 2007. Holding these
parameters fixed, I then let the preference for children fall after 2007 so
that the model matches the decline in fertility through 2023, and I solve for
the transition of the economy as prices, family composition and household
formation adjust.

I study a range of policies, starting from the 2023 point in this transition. The results show why focusing on the immediate impact of a policy can be misleading. Credit for all buyers adds births immediately, but most of these extra births disappear in four years, as many of the children born early would otherwise have been born later. Births then rise again, to about 1\% above the path without policy by the mid-2050s, as households that enter later gain more from credit while housing becomes cheaper. A supply expansion with the same effect on impact keeps raising births, to 3.0\% above the path without policy by 2063, as each cohort that enters adulthood benefits from lower housing costs. Credit for parents
raises births by 5.1\% in 2023 and by 16.6\% by 2063. In the long run,
births per household return to replacement, so the lasting gains of a policy
show up as more households rather than more children per household: with
credit for parents, the house price is 33\% higher in the new steady state,
and there are 36\% more households. A rebated property tax, by contrast,
lowers the house price by 7.0\% but raises rents by 14.5\%, and births
barely move: a lower price is not cheaper space. Housing policy aimed at
births should therefore be judged by its long-run effects, not only by its
effect on impact.

\subsection{Related work}
\label{sec:related-work}

This paper contributes to many strands of the literature on fertility, housing, and macroeconomics.

First, the paper relates to the large literature on how housing conditions shape childbearing. Studies show that fertility is lower where rent per room is higher \citep{simonTamuraHigherRents2009}, dwelling size predicts first births \citep{stromHousingFirstBirths2010}, and higher house prices raise births among owners but lower them among renters \citep{dettlingHousePricesBirth2014,lovenheimFamilyWealthShocks2013}. These papers are mostly empirical; I study these effects in a general equilibrium model, which lets me focus on complex interactions between housing, credit and fertility and their long-run patterns. Structural work has also modeled fertility and housing choices, estimating dynamic models on micro data \citep{khorunzhinaMillerAmericanDream2022,couillardBuildBabyBuild2025}. I place similar choices in a life-cycle economy where house prices respond to the number of households, and follow policies over several decades.

Second, it relates to work on credit, housing and fertility. Quasi-experimental studies find that credit raises births: households exposed to the deregulation of home lending are more likely to buy a home and have a child \citep{hacamoBabiesMortgageDeregulation2021}, lower interest rates on home loans raise births among households whose rate resets \citep{cummingMonetaryPolicyBirth2024}, federally backed home loans raised births and completed fertility during the baby boom \citep{dettlingKearney2026}, and winners of housing-credit lotteries in Brazil have more children, especially at ages 20 to 25 \citep{vandoornikFazioRamadoraiSkrastinsHousingFertility2025}. Quantitative models of housing find that easier credit mostly moves ownership rather than total housing consumption \citep{kaplanMitmanViolanteHousing2020}. My model shows that decreasing down payment requirements can raise births on impact but has a modest effect over the decades, mostly shifting the timing of births.

Third, it relates to work on earnings risk and fertility. Because a child is an irreversible commitment, prudent parents who face more risk want fewer children \citep{fraserIncomeRisk2001}. Quantitative life-cycle models find that uninsurable earnings risk delays and reduces births \citep{sommerFertilityChoiceLife2016,ejrnaesJorgensenFamilyPlanning2020,nakakuniFertilityIncomeUncertainty2026}, and \citet{deLaCroixPommeretChildbearing2021} show that risk gives the postponement of a first birth an option value. Quasi-experimental studies find that job displacement and job insecurity lower fertility \citep{delBonoWeberWinterEbmer2012,clarkLepinteurJobInsecurity2022}. I focus on modelling the housing component of such commitments. Children affect housing demand, and by modelling tenure choice and credit constraints I can better understand how the response of fertility to earnings risk is influenced by the housing market.

Fourth, it builds on quantitative models of housing over the life cycle. These models study tenure choice under earnings risk, down payments and transaction costs \citep{fisherGervaisHomeOwnership2011,halketVasudevSaving2014,boarGoreaMidriganLiquidity2022}, how credit conditions and beliefs move house prices, rents and ownership \citep{kaplanMitmanViolanteHousing2020,greenwaldGurenCredit2025}, why young households own less than they used to \citep{fisherGervaisHomeOwnership2011,pazPardoHomeownership2024}, and the intergenerational allocation of housing \citep{coven2026}. I contribute to this literature by endogenizing the fertility decision, to understand fertility and housing over the life cycle.

Finally, although this paper does not try to explain the decline in
fertility, it relates to the broad literature on its causes and
consequences. Work on the causes stresses the cost and timing of children,
bargaining within couples, changing attitudes toward parenthood, shifting priorities across cohorts and the compatibility of
career and family
\citep{jonesSchoonbroodtTertilt2010,doepkeBargainingBabiesTheory2019,kearneyPuzzleFallingUS2022,doepkeEconomicsFertilityNew2023,kimStatusExternalitiesEducation2024,kearneyLevineWhyFertilityLow2025,fernandezBerniellOnofri2026}.
Work on the consequences studies the efficiency and optimal rate of
population growth, whether low fertility stabilizes, and the economics of a
shrinking world
\citep{golosovjonestertilt2007,ballOptimalPopulation2026,coffeyGerusoSpearsVyas2026,fernandezVillaverdeNorrickTerraIncognita2026}.
I take the decline as caused by an exogenous preference shock, and study the role of housing in influencing future scenarios.

The rest of the paper is organized as follows.
Section~\ref{sec:empirical-evidence} presents the evidence on housing around
the first birth. Section~\ref{sec:model} describes the model, and
Section~\ref{sec:quant} its quantification and the transition from 2007 to
2023. Section~\ref{sec:results} describes household choices in the calibrated
model, Section~\ref{sec:policy} studies housing policies, and
Section~\ref{sec:conclusion} concludes.
